%%%% CAPÍTULO 1 - INTRODUÇÃO
%%
%% Deve apresentar uma visão global da pesquisa, incluindo: breve histórico, importância e justificativa da escolha do tema,
%% delimitações do assunto, formulação de hipóteses e objetivos da pesquisa e estrutura do trabalho.

%% Título e rótulo de capítulo (rótulos não devem conter caracteres especiais, acentuados ou cedilha)
\chapter{Introdução}\label{cap:introducao}

No âmbito da agricultura de precisão, a automação e a otimização do manejo de culturas emergiram como pilares fundamentais para garantir a segurança alimentar global e a sustentabilidade ecológica. Central a esse movimento está o Manejo Integrado de Pragas (MIP), uma abordagem estratégica que preconiza o monitoramento constante e sistemático das populações de insetos para subsidiar tomadas de decisão conscientes quanto à aplicação de defensivos agrícolas. Tradicionalmente, contudo, esse monitoramento depende da inspeção visual humana direta em campo, um processo notoriamente oneroso, lento, suscetível a falhas de amostragem e de difícil escalabilidade para propriedades de larga extensão\cite{NGUGI202127}.


Nesse cenário, a visão computacional desponta como uma alternativa tecnológica promissora, viabilizando o desenvolvimento de sistemas automatizados de diagnóstico e contagem capazes de operar de forma contínua e precisa. A popularização de técnicas de Aprendizado Profundo (Deep Learning), alavancada por arquiteturas profundas de extração de características \cite{he2016residual}, revolucionou a interpretação de imagens biológicas. Em particular, a segmentação semântica pixel a pixel baseada na arquitetura U-Net \cite{RONNEBERGER} estabeleceu novos patamares de desempenho ao unificar caminhos de contração e expansão para o mapeamento detalhado de estruturas de interesse.

Apesar dos avanços substanciais na identificação de alvos conspícuos, a aplicação prática desses modelos em ambientes agrícolas reais enfrenta barreiras complexas. Fenômenos de mimetismo cromático e camuflagem, em que a coloração, o brilho e a textura do inseto guardam profunda correlação estatística com o substrato foliar, anulam as pistas visuais estáticas predominantemente exploradas pelas Redes Neurais Convolucionais (CNNs) tradicionais \cite{Fan_2020_CVPR}. Essa limitação é severamente agravada pela natureza de pequena escala dos vetores biológicos, os quais frequentemente ocupam menos de 1\% da área total da imagem capturada. Do ponto de vista computacional, tal cenário resulta em um desbalanceamento extremo de classes durante a etapa de treinamento, induzindo as redes estáticas a convergirem para previsões triviais e errôneas, nas quais o espécime camuflado é negligenciado e assimilado à classe majoritária do fundo foliar.

Para mitigar a insuficiência dos descritores puramente espaciais fornecidos pelos canais RGB, a integração de pistas dinâmicas (temporais) surge como uma vertente teórica robusta. Fundamentada no princípio biológico de que o movimento próprio rompe o isolamento perceptivo provocado pela camuflagem, a fusão de dados baseada em modelos de duas vias (two-stream) demonstrou sucesso expressivo no reconhecimento de ações em vídeo \cite{Simonyan2014}. Ao estimar o vetor de deslocamento de pixels consecutivos por meio do algoritmo de Fluxo Óptico Denso \cite{Farneback2003}, torna-se viável codificar assinaturas cinemáticas que permanecem estritamente invariantes em relação à cor ou ao padrão textural do alvo.

A \textbf{hipótese central} deste trabalho postula que, em cenários onde a informação visual estática (domínio espacial RGB) se torna ambígua ou insuficiente, a dinâmica temporal (fluxo de movimento) atua como uma \textit{feature} discriminante invariante à camuflagem. A partir dessa premissa, levanta-se a seguinte questão de pesquisa: em que medida a incorporação explícita de informações de fluxo óptico denso pode mitigar o viés gerado pelo desbalanceamento de classes e pela ausência de contraste espectral na segmentação de insetos em pequena escala?

\section{Objetivo}

O \textbf{objetivo} geral deste Trabalho de Conclusão de Curso consiste em desenvolver, avaliar e comparar uma arquitetura híbrida de rede neural convolucional fundamentada na topologia U-Net modificada para o processamento de 6 canais (RGB concatenado às componentes cartesianas do fluxo óptico), validando a sua eficácia na segmentação semântica de insetos camuflados e de dimensões reduzidas em sequências de vídeo.

Para alcançar o objetivo geral, definem-se os seguintes \textbf{objetivos específicos}:

\begin{itemize}
    \item Curadoria e pr\'e-processamento de um \textit{dataset} anotado que abrange desde cenários de alto contraste (controle: \textit{Red/Black}) até condições de camuflagem crı́tica (desafio: \textit{Green});

    \item Desenvolvimento de um \textit{pipeline} de supervis\~ao fraca (\textit{weak supervision}) para a geração automática de \textit{Ground Truth} através de algoritmos de subtração de fundo baseados em Mistura de Gaussianas \cite{Zivkovic2004} combinados com equaliza\c{c}\~ao adaptativa \cite{Zuiderveld1994ContrastLA};

    \item Treinamento e avaliação de um modelo de refer\^encia (baseline) U-Net puramente espacial sob condi\c{c}\~oes de entrada RGB padr\~ao e sob regime de aumento de dados (\textit{data augmentation}) extensivo;

    \item Proposta e validação de uma arquitetura U-Net hı́brida, modificada para realizar a fusão de dados espaciais e temporais (Fluxo Óptico Denso), visando mitigar o severo desbalanceamento de classes inerente ao problema e contrastando quantitativamente o seu desempenho perante os modelos est\'aticos.
\end{itemize}

\section{Organização do texto}

Este documento encontra-se estruturado da seguinte forma: 
A Seção 2 apresenta a Fundamentação Teórica, incluindo conceitos aprofundados sobre a arquitetura U-Net, o algoritmo de Fluxo Óptico de Farneback, a formulação matemática da modelagem de Mistura de Gaussianas (MOG2) e da equalização adaptativa (CLAHE), bem como as métricas de validação e os trabalhos correlatos. 
A Seção 3 descreve detalhadamente os Materiais e Métodos, abordando o pipeline de pré-processamento, a formulação da rede neural híbrida e o desenho experimental em três frentes. 
Na Seção 4 encontram-se os Resultados e Discussões, avaliando exaustivamente o desempenho dos modelos em relação aos diferentes níveis de mimetismo. 
Por fim, a Seção 5 apresenta a Conclusão do trabalho, pontuando os principais achados práticos e sugerindo perspectivas de pesquisas futuras.